\documentclass[a4paper]{article}
\usepackage[paperwidth=19.304cm, paperheight=7.62cm, margin=0.1cm]{geometry}
\usepackage{fontspec,xunicode}
\usepackage{graphicx}
\usepackage{tikz}
\usepackage[absolute]{textpos}
% \usepackage[absolute,showboxes]{textpos}
\usepackage{adjustbox}
\usepackage{calc}
\usepackage{xifthen}
\setlength{\TPHorizModule}{\paperwidth}\setlength{\TPVertModule}{\paperheight}
\setlength{\parindent}{0pt}

\setmainfont{Linux Libertine}
\setsansfont{Liberation Sans}

% Info
% 
% If you have the above two fonts (and they have the same names on
% your system), then this file should compile in new TeXLives using
% xelatex lanyard.tex
% 
% This makes liberal use of the `textpos` package to place elements on
% a grid (note the difference between parens, brackets, and curly braces):
% \begin{textbock}{width}[opt center point](x position , y position)
%   stuff
% \end{textbock}
% 
% Basic text commands include
% 
% bold: \textbf{text}
% italic: \emph{text}
% size: {\fontsize{pt size}{dist between lines} text}
% font: {\fontspec{font name} text}
% 
% At least on some systems, you can get a list of valid font names by
% running `fc-list` on the command line

% defaults if something isn't defined externally
\providecommand\fname{FirstName}
\providecommand\lname{LastName}
\providecommand\titlename{Title}
\providecommand\orgname{Organization}
\providecommand\qrcode{fig/main}

\begin{document}

%% Left side

\begin{textblock}{0.2}(0.015,0.03)
  \includegraphics[width=4cm]{\qrcode}
\end{textblock}

\begin{textblock}{0.2105}[0,1](0.03,0.90) 
  \begin{small}\textsf{This QR Code \\ Contains My Contact \\
      Information.  Scan It.}\end{small}
  % alternately, you can use 
  % {\fontspec{Liberation Sans} This QR Code \\ Contains My Contact \\ Information.  Scan It.}
  % for a one-time use font
\end{textblock}

\begin{textblock}{0.06}(0.12,0.56)
  \includegraphics[width=\textwidth, angle=30]{fig/arrow}
\end{textblock}

\begin{textblock}{0.2405}(0.24,0.03)

  % First and last name should not be allowed outside of the
  % containing box, so we check that their length (at the default font
  % size) is not wider than the box.  If not, then keep it as-is.  If
  % it is, then resize it down to fit.
  {
    \def\fnamef{{\fontsize{36}{40}\textbf{\fname}}} % size of fname here
    \newdimen{\namewidth}
    \setlength{\namewidth}{\widthof{\fnamef}}
    \ifthenelse{\lengthtest{\namewidth < \textwidth}}% check if it is too long
    {\fnamef}% keep as-is
    {\resizebox{\textwidth}{!}{\fnamef}}% resize
  }

  {
    \def\lnamef{\Large{\lname}} % change size of lname here
    \newdimen{\namewidth}
    \setlength{\namewidth}{\widthof{\lnamef}}
    \ifthenelse{\lengthtest{\namewidth < \textwidth}}{\lnamef}{\resizebox{\textwidth}{!}{\lnamef}}
  }

  \vspace{2em}

  \titlename{}
  \\
  \orgname{}

\end{textblock}

\begin{textblock}{0.2105}[0,1](0.24,0.95)
  \includegraphics[width=4cm]{fig/logo}
\end{textblock}

%% Right side

% duplicated from the left, except I added 0.5 to all the x-coordinates

\begin{textblock}{0.2}(0.515,0.03)
  \includegraphics[width=4cm]{fig/speaking}
\end{textblock}

\begin{textblock}{0.2405}(0.74,0.03)

  {
    \def\fnamef{{\fontsize{36}{40}\textbf{\fname}}}%
    \newdimen{\namewidth}%
    \setlength{\namewidth}{\widthof{\fnamef}}%
    \ifthenelse{\lengthtest{\namewidth < \textwidth}}{\fnamef}{\resizebox{\textwidth}{!}{\fnamef}}%
  }

  {
    \def\lnamef{\Large{\lname}}
    \newdimen{\namewidth}
    \setlength{\namewidth}{\widthof{\lnamef}}
    \ifthenelse{\lengthtest{\namewidth < \textwidth}}{\lnamef}{\resizebox{\textwidth}{!}{\lnamef}}
  }

  \vspace{2em}

  \titlename{}
  \\
  \orgname{}

\end{textblock}

\begin{textblock}{0.2105}[0,1](0.74,0.95)
  \includegraphics[width=4cm]{fig/logo}
\end{textblock}

% Corner markers

% top left
\begin{textblock}{0.02}[0,0](0.001,0.001)
    \begin{tikzpicture}[x=\textwidth,y=\textwidth]
      \draw(0,0) to (0,1);
      \draw(0,1) to (1,1);
    \end{tikzpicture}
\end{textblock}

% top right
\begin{textblock}{0.02}[1,0](0.999,0.001)
    \begin{tikzpicture}[x=\textwidth,y=\textwidth]
      \draw(0,0) to (0,1);
      \draw(0,1) to (-1,1);
    \end{tikzpicture}
\end{textblock}

% bottom right
\begin{textblock}{0.02}[1,1](0.999,0.999)
    \begin{tikzpicture}[x=\textwidth,y=\textwidth]
      \draw(0,0) to (0,1);
      \draw(0,0) to (-1,0);
    \end{tikzpicture}
\end{textblock}

% bottom left
\begin{textblock}{0.02}[0,1](0.001,0.999)
    \begin{tikzpicture}[x=\textwidth,y=\textwidth]
      \draw(0,0) to (0,1);
      \draw(0,0) to (1,0);
    \end{tikzpicture}
\end{textblock}

% center line
% \begin{textblock}{0.01}[0.5,0](0.5,0) 
%   \begin{tikzpicture}
%     \draw(0,0) to (0,\paperheight);
%   \end{tikzpicture}
% \end{textblock}

\end{document}
